\chapter{Enchères et clôture}\label{ch:hf:auctions-and-the-close}

À 15~h~50, la bourse publie un déséquilibre acheteur de deux cent mille actions sur une valeur. Dans les dix minutes qui suivent, quelqu'un les vendra
à l'enchère, et le prix auquel elles seront vendues est fixé par celui qui affiche le meilleur cours vendeur.
Dans le carnet de clôture de ce chapitre, le seul déséquilibre ferait monter la clôture de 14~cents sur une action à 100~\$~; un fournisseur de liquidité qui offre 120\,000 actions au prix de référence limite le mouvement à 9~cents et gagne 6.9
cents par action, si 85~\% du mouvement devait revenir dès le lendemain matin. Des ordres tardifs qui modifient le déséquilibre ramènent ce gain à 5.5 cents, sauf si le
fournisseur demande une prime.

\section{L'enchère de clôture comme événement de liquidité}

L'enchère de clôture (livre~1, chapitre~13) fixe le prix auquel les fonds indiciels, les gérants qui se réfèrent à un indice et les produits dérivés se règlent. Elle constitue ainsi, pour la
journée, le plus grand réservoir de liquidité unique, et celui dont les acheteurs et les vendeurs se soucient davantage d'obtenir le cours de clôture que du prix lui-même.
Bogousslavsky et Muravyev ont mesuré l'enchère de clôture à 7.5~\% du volume quotidien américain en 2018, contre 3.1~\% en 2010. Les cours de clôture coïncident en général avec le
cours acheteur ou le cours vendeur de la période précédant la clôture~; l'impact sur le prix était plus faible que dans la négociation en continu~; et les écarts de prix constatés à l'enchère revenaient vite et presque entièrement, en
moyenne. Le volume de l'enchère rapporté au volume intrajournalier a bondi lors de l'entrée de valeurs dans le S\&P 500, puis a augmenté durablement, ce que les auteurs interprètent comme la croissance de
l'indexation et des fonds indiciels cotés qui alimentent la clôture.

Le rôle du teneur de marché y est de fournir la contrepartie d'un déséquilibre, à un prix, puis de porter la position pendant la nuit ou de la couvrir.

\section{Lire la publication du déséquilibre}

Avant l'enchère, la bourse publie ce qu'elle ferait si elle procédait dès maintenant à l'équilibrage~: le prix de référence ou le prix indicatif, les actions appariées, le
déséquilibre et son sens (la publication du déséquilibre du livre~10, chapitre~10).

\begin{dated}{2026-09}{Le calendrier de clôture d'une bourse}\label{dat:hf:auctions-and-the-close:nasdaq}
La foire aux questions de Nasdaq sur ses croisements (telle que publiée en novembre 2025) indique qu'entre 15~h~50 et 15~h~55 elle diffuse le
prix de référence courant, les actions appariées, les actions en déséquilibre et le sens du déséquilibre toutes les 10 secondes, puis, entre 15~h~55 et 16~h~00, chaque seconde, en ajoutant
les prix d'équilibre indicatifs proche et lointain. Les ordres au marché à la clôture doivent être reçus avant 15~h~55~; les ordres à cours limité à la clôture avant 15~h~58, ceux
reçus après 15~h~55 étant réajustés s'ils sont plus agressifs que les prix de référence de 15~h~50 ou de 15~h~55~; les ordres à la clôture peuvent être modifiés ou annulés avant 15~h~50.
Les ordres au déséquilibre seul existent pour compenser un déséquilibre d'ordres au marché à la clôture.
\end{dated}

Les publications sont une prévision, non la réponse~: les ordres au marché à la clôture continuent d'arriver jusqu'à leur heure limite, et les ordres à cours limité à la clôture après celle-ci. Un
fournisseur qui dimensionne son ordre d'après une publication doit tenir compte du fait que le déséquilibre peut changer avant l'équilibrage.

\section{Compenser les déséquilibres~: fixer le prix de l'ordre}

\begin{definition}[Ordre de compensation du déséquilibre]\label{def:hf:auctions-and-the-close:offset}
Un \emph{ordre de compensation du déséquilibre}\index{ordre de compensation du déséquilibre} (imbalance-offsetting order) est un ordre d'enchère de clôture placé du côté opposé à un déséquilibre publié,
fixé en prix et en taille pour n'être exécuté que si l'enchère en a besoin~: une vente face à un déséquilibre acheteur, un achat face à un déséquilibre vendeur.
\end{definition}

\begin{definition}[Courbe d'impact d'enchère]\label{def:hf:auctions-and-the-close:impact}
Une \emph{courbe d'impact d'enchère}\index{courbe d'impact d'enchère} (auction impact curve) est le déplacement du prix d'équilibre par rapport au prix de référence, en fonction de la taille d'un
déséquilibre, le reste du carnet d'enchère étant donné~; c'est l'offre de l'enchère vue depuis le déséquilibre.
\end{definition}

Dans le carnet de clôture du chapitre (300\,000 actions appariées, ordres à cours limité à la clôture de $2\,000\,j$ actions à $j$ cents de 100~\$ de chaque côté), la courbe d'impact
est concave~: 5 cents pour 25\,000 actions, 10 pour 100\,000, 14 pour 200\,000, 20 pour 400\,000. L'offre du fournisseur ramène la clôture vers le
prix de référence~; son profit est l'écart entre la clôture et la valeur du lendemain sur chaque action qu'il vend. Si une part $\pi$ du mouvement de clôture est une information qui persiste,
et que le reste revient, le fournisseur vend à la clôture au prix $p$ et anticipe la valeur du lendemain $P_0+\pi\,\Delta(I)$, où $\Delta(I)$ est le mouvement
que le seul déséquilibre aurait provoqué (voir le \cref{lst:hf:auctions-and-the-close:provider}). La meilleure offre arbitre entre la taille et le prix~: vendre davantage
abaisse la clôture pour chaque action vendue.

\begin{omfigure}
\begin{tikzpicture}
\begin{axis}[width=11cm,height=5.4cm,xlabel={\small déséquilibre (milliers d'actions)},ylabel={\small en cents},xmin=0,xmax=420,ymin=0,ymax=22,
  tick label style={font=\footnotesize},grid=major,grid style={gray!25},table/col sep=comma,
  legend style={font=\scriptsize,at={(0.5,-0.3)},anchor=north},legend columns=3]
\addplot[black,thick,mark=*,mark size=1.2pt] table[x=imbalance,y=move]{figdata/hft/16-auctions-and-the-close/provider.csv};
\addplot[black,thin,dashed,mark=o,mark size=1.2pt] table[x=imbalance,y=close15]{figdata/hft/16-auctions-and-the-close/provider.csv};
\addplot[omThm,thick,mark=square*,mark size=1.2pt] table[x=imbalance,y=ps0]{figdata/hft/16-auctions-and-the-close/provider.csv};
\addplot[omDef,thick,mark=triangle*,mark size=1.4pt] table[x=imbalance,y=ps15]{figdata/hft/16-auctions-and-the-close/provider.csv};
\addplot[gray,thick,mark=diamond*,mark size=1.4pt] table[x=imbalance,y=ps50]{figdata/hft/16-auctions-and-the-close/provider.csv};
\legend{{mouvement de clôture seul},{mouvement de clôture, avec le fournisseur},{profit par action, tout revient},{profit par action, 15~\% persistent},{profit par action, 50~\% persistent}}
\end{axis}
\end{tikzpicture}

\omcaption{La \omterm{def:hf:auctions-and-the-close:impact}{courbe d'impact d'enchère} du carnet de clôture du chapitre (le mouvement du cours de clôture au-dessus de 100~\$ face à un déséquilibre acheteur), le mouvement lorsque le fournisseur
présente sa meilleure offre (15~\% du mouvement persiste pendant la nuit), et le profit espéré par action du fournisseur à sa meilleure offre lorsque aucun, 15~\% ou la moitié du
mouvement persiste. Équilibrage calculé avec \texttt{firm.auction}. Données~: \texttt{hf\_auction.offers}.}\label{fig:hf:auctions-and-the-close:provider}
\end{omfigure}

Le profit par action croît avec le déséquilibre, moins que proportionnellement (voir la \cref{fig:hf:auctions-and-the-close:provider})~: 4.5 cents pour 100\,000 actions,
6.9 pour 200\,000 et 10 pour 400\,000 lorsque 15~\% du mouvement persiste. Il dépend autant de l'hypothèse de retour que du déséquilibre~: pour 200\,000 actions, 8
cents si tout revient, 4 si la moitié persiste. Le retour est ce que le fournisseur doit estimer à partir de l'historique, valeur par valeur, et c'est la partie de son
modèle qui peut se tromper le jour où cela compte le plus~: un déséquilibre provoqué par une nouvelle.

\section{L'enchère et le carnet continu qui la précède}

L'enchère ne se déroule pas isolément. Les traders qui voient la publication du déséquilibre peuvent négocier dans le carnet continu avant la clôture, et les
ordres à cours limité du carnet de clôture peuvent être fixés à partir de celui-ci. Un fournisseur peut couvrir sa position avant la clôture (vendre une partie du titre dans le carnet continu, ou un
contrat à terme corrélé) ou après (le cours de clôture est alors un chiffre connu, et le marché des contrats à terme continue de coter). Les contrats à terme trade at settlement (TAS) du livre~1,
chapitre~21 permettent à un fournisseur de verrouiller une couverture au prix de règlement de l'indice, en supprimant la composante de marché de son risque de nuit à une base connue.

Les ordres tardifs constituent le principal risque du fournisseur à l'intérieur de l'enchère. Avec une incertitude de 60\,000 actions sur les ordres tardifs pour un déséquilibre publié de 200\,000, l'offre
qui était la meilleure au moment de la publication (120\,000 actions au prix de référence) rapporte 5.5 cents par action au lieu de 6.9~; la meilleure offre compte tenu de
l'incertitude conserve la taille et demande 4 cents au-dessus du prix de référence, ce qui rapporte 6.0~: une prime qui protège le fournisseur lorsque le déséquilibre se réduit et
que son offre fixerait sinon la clôture.

\section{Enchères d'ouverture et reprises de cotation}

L'enchère d'ouverture et la reprise de cotation après une suspension obéissent aux mêmes mécanismes, avec moins d'information et davantage de risque~: la nouvelle survenue pendant la nuit n'est pas encore
dans aucun prix, le déséquilibre est plus petit au regard de l'incertitude de la journée, et le prix de référence est périmé. Les fournisseurs demandent davantage, et couvrent leur position sur le contrat à terme sur indice dès que la valeur ouvre.

\section{Fiches de stratégie}

\begin{strategyfile}{Fourniture de liquidité sur déséquilibre à l'enchère de clôture}\label{strat:hf:auctions-and-the-close:close}
\sfield{Qui paie, et pourquoi} Les fonds indiciels et les gérants qui se réfèrent à un indice et doivent négocier à la clôture, et qui acquittent l'impact de l'enchère pour y parvenir.
\sfield{Instruments et plateformes} Enchères de clôture des bourses de cotation~; ordres au déséquilibre seul ou ordres à cours limité à la clôture.
\sfield{Signal} Le déséquilibre publié et le prix de référence, le carnet de clôture du titre, la part estimée du mouvement qui persiste pendant la nuit.
\sfield{Dimensionnement et exécution} Offrir la taille et le prix qui maximisent le profit espéré face à la courbe d'impact~; ajouter une prime pour les ordres tardifs~; couvrir
la composante de marché par des contrats à terme sur indice au prix de règlement.
\sfield{Coûts} Risque de nuit sur la part qui persiste~; frais de la bourse~; coûts de couverture.
\sfield{Comment elle meurt} Information~: un déséquilibre provoqué par une nouvelle laisse l'essentiel du mouvement en place (4 cents par action lorsque la moitié persiste, au lieu de 6.9)~;
concurrence entre fournisseurs.
\sfield{Horizon, capacité, infrastructure} Dix minutes dans la phase de clôture et jusqu'au lendemain~; capacité issue des déséquilibres de milliers de valeurs~; les
flux de déséquilibre de chaque bourse.
\sfield{Backtest honnête} Publications telles que diffusées à chaque instant, ordres tardifs, règles d'équilibrage, et ouverture du lendemain comme sortie.
\sfield{Sources} Bogousslavsky et Muravyev (2023)~; l'encadré daté.
\end{strategyfile}

\begin{strategyfile}{Opening-auction liquidity provision}\label{strat:hf:auctions-and-the-close:open}
\sfield{Qui paie le fournisseur, et pourquoi} Des traders qui doivent exécuter à l'ouverture après des nouvelles ou des flux de la nuit.
\sfield{Instruments et plateformes} Enchères d'ouverture~; publications de déséquilibre avant l'ouverture.
\sfield{Signal} Le déséquilibre d'ouverture et le cours indicatif, comparés à la fair value issue des contrats à terme de la nuit et des nouvelles.
\sfield{Taille et exécution} Des tailles plus petites et des primes plus larges qu'à la clôture~; couvrir sur le contrat à terme dès que l'action ouvre.
\sfield{Coûts} Information de la nuit absente du prix de référence~; volatilité des premières minutes.
\sfield{Comment elle meurt} Les nouvelles que la fair value du fournisseur manque.
\sfield{Horizon, capacité, infrastructure} Quelques minutes~; les données de pré-ouverture de chaque plateforme.
\sfield{Backtest honnête} Contrats à terme et nouvelles tels qu'ils sont connus à chaque heure de pré-ouverture.
\sfield{Sources} Livre~1, chapitre~13~; le modèle d'enchère de clôture de ce chapitre, appliqué avec moins d'information.
\end{strategyfile}

\begin{strategyfile}{Clôture couverte par des contrats à terme trade-at-settlement}\label{strat:hf:auctions-and-the-close:tas}
\sfield{Qui paie le fournisseur, et pourquoi} Comme pour la fourniture de liquidité à la clôture, la part de marché du risque de la nuit étant transférée à une base connue.
\sfield{Instruments et plateformes} Enchères de clôture des composantes de l'indice~; contrats à terme trade-at-settlement sur indice (livre~1, chapitre~21).
\sfield{Signal} Déséquilibres agrégés sur les composantes de l'indice, qui prédisent le mouvement de clôture de l'indice.
\sfield{Taille et exécution} Fournir de la liquidité face aux déséquilibres des composantes~; acheter ou vendre des contrats à terme trade-at-settlement pour le bêta agrégé au cours de la
journée, avant que la clôture soit connue.
\sfield{Coûts} La base trade-at-settlement, les frais, la part idiosyncratique laissée non couverte.
\sfield{Comment elle meurt} L'encombrement (crowding) sur la base trade-at-settlement~; des déséquilibres qui se compensent d'un titre à l'autre, de sorte que la couverture dépasse la cible.
\sfield{Horizon, capacité, infrastructure} La dernière heure et la nuit.
\sfield{Backtest honnête} Les prix trade-at-settlement disponibles à chaque instant, et non le règlement lui-même.
\sfield{Sources} Livre~1, chapitre~21.
\end{strategyfile}

\section{Tutoriel~: deux cent mille actions à vendre}\label{tut:hf:auctions-and-the-close:tut}

\begin{tutorial}
\textbf{Objectif.} Tarifer un \omterm{def:hf:auctions-and-the-close:offset}{ordre de compensation du déséquilibre} à partir de la \omterm{def:hf:auctions-and-the-close:impact}{courbe d'impact d'enchère} et du retournement overnight, et mesurer le coût des ordres tardifs.
\textbf{État final~:} \cref{fig:hf:auctions-and-the-close:provider} et les chiffres du texte.
\begin{tutsteps}
\item \textbf{Le carnet de clôture} et ses ordres pour \texttt{firm.auction.uncross}.
\omcode{code/firm/auctionmm/firm_auctionmm.py}{35}{50}{Intérêt pairé au marché à la clôture, ordres à cours limité à la clôture s'épaississant à distance de la référence, déséquilibre et fournisseur.}\label{lst:hf:auctions-and-the-close:book}
\item \textbf{Le P\&L du fournisseur} face à la valeur du lendemain, avec les ordres tardifs.
\omcode{code/firm/auctionmm/firm_auctionmm.py}{63}{80}{Croisement avec et sans fournisseur~; la valeur du lendemain conserve une part du mouvement que le seul déséquilibre aurait causé.}\label{lst:hf:auctions-and-the-close:provider}
\item \textbf{Le meilleur cours vendeur} sur des tailles de 20\,000 à 400\,000 actions et des limites de 0 à 10 cents au-delà de la référence (\texttt{hf\_auction.offers}).
\item \textbf{Ordres tardifs}~: un écart type de 60\,000 actions sur un déséquilibre publié de 200\,000 (\texttt{hf\_auction.late}).
\end{tutsteps}
\textbf{Pour aller plus loin.} Ajouter un second fournisseur et chercher l'équilibre~; estimer la part qui subsiste à partir de l'historique
des écarts de clôture de chaque titre~; lancer l'enchère sur \texttt{firm.exchsim}, avec le preset \texttt{auction\_venue} et ses publications.
\end{tutorial}

\section{Construction~: le module de liquidité des enchères}\label{bld:hf:auctions-and-the-close:auctionmm}

\begin{build}
\bfield{Rôle} Modéliser les carnets de clôture, les courbes d'impact et les \omterm{def:hf:auctions-and-the-close:offset}{ordres de compensation du déséquilibre}, avec les ordres tardifs et le retournement overnight.
\bfield{Interface} \texttt{ClosingBook(p0, levels, d0, paired)} avec \texttt{orders(imbalance, q, k)}, \texttt{close\_price}, \texttt{impact\_curve},
\texttt{provider(book, imbalance, q, k, perm, late\_sd, n, seed)}, \texttt{best\_offer}, \texttt{publications}. Construit sur \texttt{firm.auction}
(livre~1).
\bfield{Règles} Prix en ticks~; l'ordre du fournisseur est un ordre à cours limité situé à $k$ ticks au-delà de la référence, du côté opposé au déséquilibre~; P\&L face
à la valeur du lendemain.
\bfield{Tests d'acceptation} \texttt{code/firm/auctionmm/tests/}~: un carnet équilibré clôture à la référence~; un déséquilibre de 30 dans un carnet jouet clôture à deux
ticks de part et d'autre~; la courbe d'impact croît~; le fournisseur déplace la clôture et s'exécute~; son P\&L baisse quand le mouvement subsiste~; le meilleur cours vendeur est
au moins aussi bon qu'un cours donné~; les publications convergent vers le déséquilibre final.
\bfield{Pour aller plus loin} Plusieurs fournisseurs~; couverture avec des contrats à terme trade-at-settlement~; enchères d'ouverture avec l'information de la nuit.
\end{build}

\begin{omsources}
\item V.~Bogousslavsky, D.~Muravyev, Who trades at the close? Implications for price discovery and liquidity, \emph{Journal of Financial
Markets} 66, 2023, 100852.
\item Nasdaq, The Nasdaq Opening and Closing Crosses: frequently asked questions (novembre 2025).
\end{omsources}

\section{Exercices}

\begin{exercise}[$\star$]\label{exo:hf:auctions-and-the-close:1}
Dans un carnet de clôture jouet, des ordres de vente à cours limité à la clôture reposent à 101 (10 actions), 102 (20) et 103 (30) pour une référence de 100. À quel
prix l'enchère clôture-t-elle un déséquilibre d'achat de 30 actions au marché à la clôture~?
\end{exercise}

\begin{exercise}[$\star$]\label{exo:hf:auctions-and-the-close:2}
La clôture se déplace de 9 cents avec l'ordre du fournisseur~; la valeur du lendemain conserve 15~\% du mouvement de 14 cents que le seul déséquilibre aurait causé.
Calculer la part de ce que gagne le fournisseur par action, puis son gain sur 120\,000 actions.
\end{exercise}

\begin{exercise}[$\star$]\label{exo:hf:auctions-and-the-close:3}
Selon le calendrier de l'encadré daté, jusqu'à quelle heure un ordre au marché à la clôture peut-il être saisi, et à partir de quand le déséquilibre commence-t-il à être publié toutes les secondes~?
\end{exercise}

\begin{exercise}[$\star\star$]\label{exo:hf:auctions-and-the-close:4}
Pourquoi le meilleur cours vendeur du fournisseur remplit-il rarement la totalité du déséquilibre~?
\end{exercise}

\begin{exercise}[$\star\star$]\label{exo:hf:auctions-and-the-close:5}
Pourquoi le profit par action tombe-t-il de 8 à 4 cents pour 200\,000 actions lorsque la moitié du mouvement subsiste au lieu de zéro~?
\end{exercise}

\begin{exercise}[$\star\star$]\label{exo:hf:auctions-and-the-close:6}
Pourquoi une prime de 4 cents aide-t-elle lorsque les ordres tardifs rendent le déséquilibre incertain~?
\end{exercise}

\begin{exercise}[$\star\star\star$]\label{exo:hf:auctions-and-the-close:7}
\emph{Programmation.} Doubler la profondeur du carnet de clôture (\texttt{d0=4000}). Que se passe-t-il pour l'impact d'un déséquilibre de 200\,000 actions et pour le
meilleur profit par action du fournisseur avec 15~\% qui subsistent~?
\end{exercise}

\begin{exercise}[$\star\star\star$]\label{exo:hf:auctions-and-the-close:8}
\emph{Trouver la faille.} «~Les écarts de clôture reviennent presque entièrement en moyenne, donc nous compensons chaque déséquilibre à taille pleine.~»
\end{exercise}

\section{Problème~: deux cent mille actions à vendre}

\begin{problem}[{Problème du week-end --- deux cent mille actions à vendre}]\label{pb:hf:auctions-and-the-close:1}
Un fournisseur décide à 15 h 55 combien d'actions offrir face à un déséquilibre acheteur publié de 200\,000.

\medskip\noindent\textbf{Partie I --- L'enchère.}
\begin{enumerate}
\item Pourquoi la clôture est-elle le plus grand réservoir de liquidité de la journée~?
\item Que Bogousslavsky et Muravyev ont-ils établi sur les volumes, les prix et les retournements à la clôture~?
\item Résumer le calendrier de l'encadré daté.
\item Que contient une publication de déséquilibre~?
\end{enumerate}

\medskip\noindent\textbf{Partie II --- Tarification.}
\begin{enumerate}[resume]
\item Définir un \omterm{def:hf:auctions-and-the-close:offset}{ordre de compensation du déséquilibre}.
\item Définir la \omterm{def:hf:auctions-and-the-close:impact}{courbe d'impact d'enchère} et donner ses valeurs dans le livre du chapitre.
\item Écrire le profit du fournisseur étant donnée la part $\pi$ qui reste.
\item Pourquoi vendre davantage réduit-il le profit par action~?
\end{enumerate}

\medskip\noindent\textbf{Partie III --- Mesures.}
\begin{enumerate}[resume]
\item Donner le meilleur cours vendeur et le profit pour 200\,000 actions, avec 15~\% qui restent.
\item Comment le profit par action varie-t-il avec le déséquilibre et avec $\pi$~?
\item Que coûtent les ordres tardifs, et qu'est-ce qui en récupère une partie~?
\item De combien le fournisseur déplace-t-il la clôture~?
\end{enumerate}

\medskip\noindent\textbf{Partie IV --- Le verdict.}
\begin{enumerate}[resume]
\item Énoncer le \emph{résultat nommé}~: le profit espéré par action du fournisseur de compensation face à la taille du déséquilibre, et la part du
mouvement de clôture qui revient pendant la nuit.
\item Comment estimer $\pi$ action par action~?
\item Comment couvrir la position de nuit~?
\item Qu'est-ce qui change à l'ouverture~?
\item Que fait un second fournisseur au profit du premier~?
\item Quelle contribution le fournisseur apporte-t-il à la qualité du cours de clôture~?
\item Quelle fiche de stratégie est la plus exposée aux nouvelles~?
\item En une phrase~: que vend un fournisseur de déséquilibre~?
\end{enumerate}
\end{problem}

\section{Questions d'entretien}

\begin{interviewq}[$\star$ \iqroles{trader}]\label{iq:hf:auctions-and-the-close:1}
Une action affiche un déséquilibre acheteur de 500\,000 actions à 15 h 50. Qu'est-ce que vous regardez avant de proposer une contrepartie~?
\end{interviewq}

\begin{interviewq}[$\star\star$ \iqroles{chercheur}]\label{iq:hf:auctions-and-the-close:2}
Comment mesureriez-vous la part des mouvements de cours à la clôture de l'enchère qui revient pendant la nuit~?
\end{interviewq}

\begin{interviewq}[$\star\star$ \iqroles{développeur}]\label{iq:hf:auctions-and-the-close:3}
Implémentez la règle de détermination du prix d'équilibre d'une enchère. Comment traitez-vous les égalités~?
\end{interviewq}

\begin{interviewq}[$\star\star$ \iqroles{risque}]\label{iq:hf:auctions-and-the-close:4}
Votre desk compense chaque jour des déséquilibres sur 800 actions. Quel est le risque de nuit de la société, et comment est-il limité~?
\end{interviewq}

\begin{interviewq}[$\star\star$ \iqroles{trader}]\label{iq:hf:auctions-and-the-close:5}
Pourquoi les cours de clôture peuvent-ils être déformés les jours de rééquilibrage d'indice, et qui en profite~?
\end{interviewq}

\begin{interviewq}[$\star\star\star$ \iqroles{chercheur}]\label{iq:hf:auctions-and-the-close:6}
Avec une courbe d'impact linéaire $\Delta(I)=\lambda I$, une part $\pi$ qui reste et un fournisseur qui propose $q$ au cours de référence, trouvez le $q$ optimal.
\end{interviewq}
